\section{The negative binomial distribution}

We say that $X$ has the negative binomial
distribution with parameters $m\ge1$ and $0<p<1$, and
denote $X\sim \operatorname{NB}(m,p)$, if $X$
satisfies

\[
\mathrm{P}(X=k) = \binom{k+m-1}{k} p^m (1-p)^k \qquad (k\ge 0).
\]

\noindent $\textbf{Properties:}$

\begin{enumerate}[label=\arabic*.]
\item $\operatorname{NB}(1,p) \stackrel{d}{=}
\operatorname{Geom_0}(1-p)$.

\item $\operatorname{NB}(m,p) \boxplus
\operatorname{NB}(n,p) = \operatorname{NB}(n+m,p)$.

\item If $W_1,W_2,W_3, \ldots$ is a sequence of i.i.d.
r.v.'s with distribution $\operatorname{Ber}(p)$,
then

\[
X = \min\left\{ k\ge 0\,:\, \sum_{j=1}^{k+m} (1-W_j) = m\right\}  \sim \operatorname{NB}(m,p).
\]

\noindent In words, when performing a sequence of identical
experiments, each with probability $p$ of success,
the number of successes observed before the $m$th
failure is distributed according to
$\operatorname{NB}(m,p)$.

\item $\operatorname{NB}(m,p) =
\mathop{\boxplus}\limits_{k=1}^m
\operatorname{Geom_0}(1-p)$.
\end{enumerate}
